% Tapis de Quillen, with Quillen's offprint and preprint (cote 162-5, typed part): archive cote 162-5, site page 11 (PDF page 12).
% typescript. Draft compared: none.
% AI-transcribed by Claude Opus 5.5 (claude-opus-5-5, adaptive thinking, effort high), Message Batches API, batch t162_full (msgbatch_01UAh3zyWPd454AeDkhnmQ33), 2026-11-13; not yet checked by eye.
\scanpage{162-5/11}
\mspage{10}

qu'i\unsure{l} conviendra de poser les définitions adéquates). - Peut-être Quill
\margin{\unsure{sic}}a-t-il une définition semi-simpliciale plus directe, sans passer par des
n-catégories, des $K_i$ et plus précisément du H-espace qui devrait correspo
\inserted[typed]{en général}
dre à C. \xout{\illegible{2}} Ce n'est certainement pas l'ensemble semi-simplicial
SC \xout{\unsure{engendrés}} associé à C, puisque dans le cas d'une catégorie additive,
C est homotopiquement trivial (en effet le foncteur identique est ``homoto
au foncteur nul, puisqu'il suffit d'avoir un homomorphisme d'un foncteur
dans l'autre; d'où une équivalence d'homotopie entre C et la catégorie
nulle). Pourtant, il semble bien que dans le cas où C est une Gr-catégor
\margin{\unsure{?}}i.e. admet des inverses, alors ce soit bien la construction ``naïve'' SC qu
est la bonne, puisqu'on trouve bien un $\hand{\pi}_0$ et un $\hand{\pi}_1$ qui sont aussi les
$\hand{\pi}_0$ et $\hand{\pi}_1$ au sens catégorie sans Gr-structure.

\hand{\unsure{À expliquer pour} des Gr-n-catégories, les $\pi_i$ au sens}
\margin{\unsure{?}}\hand{semi-simplicial sont aussi les $\pi_i$ au sens Gr-n-catégories,}
\hand{savoir les objets à isom près dans divers \uline{Hom}$(1,1)$\,..... Il y a}
\hand{aussi une interprétation \unsure{de l'opération} \deleted{\unsure{Steenrod} \illegible{2}}}
\hand{\deleted{\unsure{constructions d'opérations des shifts géométriques évidentes}} sur les}
\hand{n-catégories.....}

\margin{11}
